\documentclass[12pt,a4paper]{article}
\usepackage{amsmath,amssymb,amsthm,hyperref,geometry}
\usepackage{enumitem,booktabs,array}
\usepackage[T1]{fontenc}
\usepackage[utf8]{inputenc}
\usepackage[ngerman]{babel}
\geometry{margin=2.5cm}

% -------------------------------------------------------
% Theorem-Umgebungen
% -------------------------------------------------------
\newtheorem{theorem}{Satz}[section]
\newtheorem{lemma}[theorem]{Lemma}
\newtheorem{corollary}[theorem]{Korollar}
\newtheorem{proposition}[theorem]{Proposition}
\newtheorem{conjecture}[theorem]{Vermutung}
\newtheorem{definition}[theorem]{Definition}

\theoremstyle{remark}
\newtheorem{remark}[theorem]{Bemerkung}
\newtheorem{example}[theorem]{Beispiel}

% -------------------------------------------------------
% Metadaten
% -------------------------------------------------------
\title{\textbf{Die $abc$-Vermutung}\\[0.5em]
\large Radikal, Höhenfunktion, Mochizukis IUT-Theorie und Konsequenzen}
\author{Michael Fuhrmann}
\date{März 2026}

\hypersetup{
  pdftitle={Die abc-Vermutung},
  pdfauthor={Michael Fuhrmann},
  colorlinks=true,
  linkcolor=blue,
  citecolor=blue
}

\begin{document}
\maketitle
\begin{abstract}
Die $abc$-Vermutung, unabhängig von Oesterlé und Masser im Jahr 1985 aufgestellt, ist eines
der weitreichendsten offenen Probleme der Zahlentheorie. Sie postuliert eine tiefgreifende
Ungleichung zwischen den Größen dreier teilerfremder ganzer Zahlen $a, b, c$ mit $a+b=c$
und dem Produkt ihrer verschiedenen Primfaktoren — dem sogenannten \emph{Radikal}
$\operatorname{rad}(abc)$. Die Vermutung impliziert als Spezialfälle den Großen Fermatschen
Satz für große Exponenten, den Satz von Roth, Szpiros Vermutung über elliptische Kurven
und zahlreiche klassische diophantische Resultate.

Im Jahr 2012 kündigte Shinichi Mochizuki einen Beweis mittels seiner
\emph{Inter-Universellen Teichmüller-Theorie} (IUT) an. Stand 2026 ist die mathematische
Gemeinschaft gespalten: Der Einwand von Scholze und Stix (2018) bleibt ungelöst.
Dieses Paper gibt einen Überblick über die Formulierung der Vermutung, ihre wichtigsten
Konsequenzen, die Strategie von Mochizukis Ansatz und den aktuellen Stand der Kontroverse.
\end{abstract}

\tableofcontents
\newpage

% ================================================================
\section{Formulierung der Vermutung}
% ================================================================

\subsection{Das Radikal}

\begin{definition}[Radikal]
  Für eine positive ganze Zahl $n$ ist das \emph{Radikal} von $n$ definiert als
  \[
    \operatorname{rad}(n) \;:=\; \prod_{\substack{p\;\text{prim} \\ p \mid n}} p.
  \]
  Mit anderen Worten: $\operatorname{rad}(n)$ ist der größte quadratfreie Teiler von $n$.
\end{definition}

\begin{example}
  $\operatorname{rad}(12) = \operatorname{rad}(2^2\cdot 3) = 2\cdot 3 = 6$;
  $\quad \operatorname{rad}(360) = \operatorname{rad}(2^3\cdot 3^2\cdot 5) = 30$.
\end{example}

\subsection{Die Höhenfunktion}

Die \emph{logarithmische Höhe} (oder \emph{Qualität}) eines Tripels $(a, b, c)$ mit
$a+b=c$, $\ggT(a,b)=1$, $a,b,c>0$ ist definiert als
\begin{equation}\label{eq:qualitaet}
  q(a,b,c) \;:=\; \frac{\log\max(a,b,c)}{\log\operatorname{rad}(abc)}
             \;=\; \frac{\log c}{\log\operatorname{rad}(abc)}.
\end{equation}
Die Höhenfunktion $h(a,b,c) = \log\max(|a|,|b|,|c|)$ ist die additive Entsprechung dazu.
Die $abc$-Vermutung besagt im Kern, dass die Qualität $q(a,b,c)$ nicht wesentlich
größer als $1$ sein kann.

\subsection{Die $abc$-Vermutung}

\begin{conjecture}[$abc$-Vermutung, Oesterlé--Masser 1985~\cite{Oesterle1988,Masser1985}]
  Für jedes $\varepsilon > 0$ existiert eine Konstante $K_\varepsilon > 0$, sodass für
  alle Tripel $(a,b,c)$ teilerfremder positiver ganzer Zahlen mit $a + b = c$ gilt:
  \begin{equation}\label{eq:abc}
    c \;\leq\; K_\varepsilon \cdot \operatorname{rad}(abc)^{1+\varepsilon}.
  \end{equation}
\end{conjecture}

Äquivalent dazu gilt: Nur endlich viele Tripel haben die Qualität $q(a,b,c) > 1+\varepsilon$.

\begin{remark}
  Der Term $\varepsilon$ ist notwendig: Das Rekordtripel $a=1$, $b=3^{10}\cdot109$,
  $c=23^5$ hat Qualität $q \approx 1{,}6299 > 1$, also $c = \operatorname{rad}(abc)^{1{,}6299}
  \gg \operatorname{rad}(abc)$, was zeigt, dass die Schranke $c \leq K \cdot \operatorname{rad}(abc)$
  (Exponent $1$ ohne $\varepsilon$) nicht universell gelten kann.  Tripel mit Qualität
  $q \to 1$ von unten existieren in großer Zahl (z.\,B.\ $a=1$, $b=2^n-1$, $c=2^n$,
  wenn $2^n-1$ eine Mersenne-Primzahl ist), verletzen die $\varepsilon=0$-Schranke
  jedoch \emph{nicht} — es sind die Rekordtripel mit $q > 1$, die $\varepsilon > 0$
  erzwingen.  Die größte bekannte Qualität oberhalb von $1$ beträgt $q \approx 1{,}6299$.
\end{remark}

\subsection{Äquivalente Formulierungen}

Eine äquivalente Form nutzt die Höhe $h(a,b,c) = \log\max(|a|,|b|,|c|)$: Die
$abc$-Vermutung besagt dann
\[
  h(a,b,c) \;\leq\; (1+\varepsilon)\sum_{p\mid abc} \log p + O_\varepsilon(1)
\]
für alle teilerfremden $(a,b,c)$ mit $a+b+c=0$ (projektive Formulierung mit Vorzeichen).

Eine schwache Form postuliert die Existenz von $K$ (abhängig von $\varepsilon$), sodass
für \emph{alle} teilerfremden Tripel gilt:
\[
  \max(|a|,|b|,|c|) \;\leq\; K \cdot \operatorname{rad}(abc)^2.
\]
Diese schwache $abc$-Vermutung (Exponent $2$ statt $1+\varepsilon$) folgt aus der Theorie
der linearen Formen in Logarithmen (Bakers Satz) und ist damit unbedingt wahr.

% ================================================================
\section{Hochqualitative $abc$-Tripel}
% ================================================================

Systematische Computersuchen (Reyssat, Nitaj,~\cite{Nitaj1996}) haben Tripel mit großer
Qualität katalogisiert. Die zehn besten bekannten Tripel haben alle $q > 1{,}4$:

\begin{table}[h]
\centering
\begin{tabular}{lll}
\toprule
$a$ & $b$ & $c = a+b$ \quad (Qualität) \\
\midrule
$1$ & $2$ & $3$ \quad ($q=1$, trivial) \\
$1$ & $8$ & $9 = 3^2$ \quad ($q = \log 9/\log 6 \approx 1{,}226$) \\
$2^5$ & $7^2$ & $3^4$ \quad ($q\approx1{,}292$, alle klein) \\
$1$ & $3^{10}\!\cdot\!109$ & $23^5$ \quad ($q\approx1{,}630$, aktueller Rekord) \\
\bottomrule
\end{tabular}
\caption{Ausgewählte hochqualitative $abc$-Tripel}
\end{table}

Die Goldfeld-Vermutung, eine Verfeinerung, sagt $N(Q,X) \sim C_\varepsilon X^{2-2/Q}$
für die Anzahl der Tripel mit $c\leq X$ und Qualität $>Q$ voraus.

% ================================================================
\section{Konsequenzen der $abc$-Vermutung}
% ================================================================

\subsection{Der Große Fermatsche Satz für große Exponenten}

\begin{theorem}[$abc$ $\Rightarrow$ FLT für $n \geq 6$~\cite{Granville1990}]
  Unter Annahme der $abc$-Vermutung hat für alle $n \geq 6$ die Gleichung
  $x^n + y^n = z^n$ keine Lösung in positiven ganzen Zahlen.
\end{theorem}

\begin{proof}[Beweis]
  Angenommen $x^n + y^n = z^n$ mit $\ggT(x,y)=1$. Setze $a=x^n$, $b=y^n$, $c=z^n$.
  Dann gilt $\operatorname{rad}(abc) \leq xyz \leq z^3$. Die $abc$-Vermutung (mit
  $\varepsilon=1$) liefert
  \[
    c = z^n \;\leq\; K_1 \cdot (z^3)^2 = K_1 z^6,
  \]
  also $z^{n-6}\leq K_1$, was $z$ (und damit $x, y$) für $n>6$ beschränkt. Eine
  Verfeinerung mit $\varepsilon < 1$ erweitert das Argument auf $n\geq 6$.
\end{proof}

\begin{remark}
  Wiles' Beweis des FLT (1995) gilt für alle $n\geq 3$ und ist unbedingt gültig.
  Der $abc$-basierte Beweis ist zwar viel einfacher, setzt aber die unbewiesene
  $abc$-Vermutung voraus und deckt zunächst nur $n\geq 6$ ab.
\end{remark}

\subsection{Roths Satz und Thue--Mahler-Gleichungen}

\begin{corollary}[$abc$ $\Rightarrow$ Satz von Roth]
  Die $abc$-Vermutung impliziert: Für jede algebraische Zahl $\alpha$ vom Grad $d\geq 2$
  und jedes $\varepsilon>0$ gibt es nur endlich viele rationale Zahlen $p/q$ mit
  $|\alpha - p/q| < q^{-2-\varepsilon}$.
\end{corollary}

Genauer genommen impliziert $abc$ die volle Stärke von Bakers Theorie der linearen Formen
in Logarithmen (und darüber hinaus), wodurch alle effektiven diophantischen
Approximationsresultate einheitlich werden.

\subsection{Szpiros Vermutung für elliptische Kurven}

Sei $E/\mathbb{Q}$ eine elliptische Kurve mit minimaler Diskriminante $\Delta_{\min}$
und Führer $N$.

\begin{conjecture}[Szpiros Vermutung~\cite{Szpiro1981}]
  Für jedes $\varepsilon>0$ existiert $C_\varepsilon$, sodass
  \begin{equation}\label{eq:szpiro}
    |\Delta_{\min}| \;\leq\; C_\varepsilon \cdot N^{6+\varepsilon}.
  \end{equation}
\end{conjecture}

\begin{theorem}[Frey--Masser--Oesterlé; siehe~\cite{Silverman1994}]
  Die $abc$-Vermutung ist äquivalent zu Szpiros Vermutung \eqref{eq:szpiro} für
  halbstabile elliptische Kurven.
\end{theorem}

Die Verbindung stellt die \emph{Frey-Kurve} her: Zu einem $abc$-Tripel hat die
elliptische Kurve $y^2 = x(x-a^n)(x+b^n)$ eine Diskriminante, die mit $(abc)^{2n}$
zusammenhängt, und einen Führer, der im Wesentlichen $\operatorname{rad}(abc)^2$ ist.

\subsection{Weitere Konsequenzen}

Die $abc$-Vermutung impliziert oder vereinfacht Beweise von:
\begin{itemize}[nosep]
  \item \textbf{Mordell--Faltings:} Die Faltings-artige Schranke $h(P)\leq c(E)$
    für Höhen rationaler Punkte.
  \item \textbf{Wieferich-Primzahlen:} Nur endlich viele Wieferich-Primzahlen
    (Primzahlen $p$ mit $2^{p-1}\equiv1\pmod{p^2}$) — bedingt auf $abc$
    (Granville--Langevin 1985).
  \item \textbf{Halls Vermutung:} $|x^3 - y^2| \gg x^{1/2 - \varepsilon}$
    für verschiedene ganze Zahlen, folgt aus $abc$.
  \item \textbf{Catalans Vermutung:} Nur $2^3$ und $3^2$ sind aufeinanderfolgende
    Potenzen (unbedingt bewiesen von Mihailescu 2002, aber einfacher unter $abc$).
  \item \textbf{Verallgemeinerter Fermat:} $x^p + y^q = z^r$ hat endlich viele
    teilerfremde Lösungen für $1/p+1/q+1/r < 1$.
\end{itemize}

% ================================================================
\section{Mochizukis Inter-Universelle Teichmüller-Theorie (IUT)}
% ================================================================

\subsection{Überblick}

In vier Preprints, die im August 2012 veröffentlicht wurden~\cite{Mochizuki2012a,Mochizuki2012b,
Mochizuki2012c,Mochizuki2012d}, kündigte Shinichi Mochizuki (RIMS, Kyoto) einen Beweis der
$abc$-Vermutung durch ein grundlegend neues Rahmenwerk an, das er als
\emph{Inter-Universelle Teichmüller-Theorie} (IUT, auch \emph{arithmetische Deformationstheorie}
genannt) bezeichnet.

\subsection{Schlüsselideen der IUT}

IUT arbeitet mit Objekten, die als \emph{Hodge-Theater} bezeichnet werden — aufwendigen
kombinatorischen und kategorischen Strukturen, die aus Daten einer elliptischen Kurve über
einem Zahlkörper aufgebaut sind. Das zentrale Konstrukt umfasst:

\begin{enumerate}[nosep]
  \item \textbf{Theater und Verbindungen:} Ein \emph{$\Theta^{\pm\text{ell}}$NF-Hodge-Theater}
    bündelt die Arithmetik der elliptischen Kurve und einer Hilfsprimzahl $l$ in einem
    kategorischen Objekt. \emph{Étale Theta-Funktionen} kodieren die $p$-adische Geometrie.
  \item \textbf{Multiradiality:} Mochizuki führt eine Unterscheidung zwischen
    \emph{mono-anabelschen} und \emph{poly-anabelschen} Strukturen ein. Das Schlüsselprinzip
    ist, dass bestimmte Daten allein aus der absoluten Galois-Gruppe rekonstruiert werden
    können (\emph{mono-anabelsche Rekonstruktion}).
  \item \textbf{Log-Theta-Gitter:} Aufeinanderfolgende Anwendungen des \emph{Log-Links}
    (der von multiplikativen zu additiven Strukturen via $p$-adische Logarithmen übergeht)
    und des \emph{Theta-Links} (der verschiedene Theater miteinander verbindet) bilden
    ein Gitter.
  \item \textbf{Kummer-Theorie und Unbestimmtheiten:} Drei Typen von Unbestimmtheit
    (bezeichnet als $\mathrm{Ind}1$, $\mathrm{Ind}2$, $\mathrm{Ind}3$) entstehen aus der
    Nicht-Kommutativität der Log- und Theta-Links. Die Kontrolle dieser Unbestimmtheiten
    durch explizite Schranken liefert die $abc$-Ungleichung.
\end{enumerate}

Die Beweisstrategie besteht darin zu zeigen:
\[
  -\log|\Delta_E| \;\leq\; (1+\varepsilon)\log\operatorname{rad}(\Delta_E) + O(1),
\]
was (via Szpiros Äquivalenz) in die $abc$-Vermutung übersetzt wird.

\subsection{Der Einwand von Scholze und Stix (2018)}

Im März 2018 besuchten Peter Scholze und Jakob Stix Mochizuki in Kyoto und verfassten
anschließend einen kritischen Bericht~\cite{ScholzeStix2018}, in dem sie einen
fundamentalen Fehler in der IUT identifiziert zu haben behaupten.

Ihr zentraler Einwand betrifft Mochizukis Behauptung, dass bestimmte Identifikationen
von Ringstrukturen (die er als ``Theta-Link'' bezeichnet) mit der Ringstruktur auf beiden
Seiten kompatibel gemacht werden können. Scholze und Stix argumentieren, dass die
relevante Ungleichung nach dem Theta-Link auf eine triviale Identität $\log(q) \leq \log(q)$
reduziert wird und daher keinen Informationsgehalt trägt. Ihrer Ansicht nach sind die
Unbestimmtheiten zu grob, um eine nicht-triviale Schranke zu etablieren.

Mochizuki antwortete mit ausführlichen Gegenargumenten~\cite{Mochizuki2018rebuttal} und
behauptete, Scholze und Stix hätten die Rollen der verschiedenen ``Kopien'' des Rings
falsch identifiziert.

\subsection{Aktuelle Situation (2026)}

Die japanische Zeitschrift \textit{Publications of RIMS} veröffentlichte Mochizukis vier
IUT-Papers im März 2021 nach einem internen Begutachtungsprozess. Dies hat die Kontroverse
jedoch nicht gelöst:
\begin{itemize}[nosep]
  \item Der Einwand von Scholze und Stix wird von vielen führenden Zahlentheoretikern
    nach wie vor als stichhaltig angesehen (Scholze, Stix, Conrad, de Jong, \ldots).
  \item Eine kleine Gruppe von Forschern (Fesenko, Yamashita, Minamide, \ldots) hat die
    IUT eingehend studiert und bestätigt deren Korrektheit.
  \item Eine unabhängige Verifikation durch breiten Konsens ist bisher nicht zustande
    gekommen.
  \item Eine vereinfachte Darstellung der IUT steht weiterhin aus.
\end{itemize}

Die mathematische Gemeinschaft befindet sich effektiv in einer Sackgasse, und die
$abc$-Vermutung gilt im Mainstream-Konsens als nicht bewiesen.

% ================================================================
\section{Alternative Ansätze}
% ================================================================

\subsection{$p$-adische und Arakelov-Geometrie}

Mehrere Forscher haben untersucht, ob Elemente der IUT durch konventionellere $p$-adische
Hodge-Theorie oder Arakelov-Schnitttheorie ersetzt werden können. Tan~\cite{Tan2022} und
andere haben partielle Umformulierungen skizziert; bisher gelang es niemandem, den
Theta-Link zu umgehen.

\subsection{Funktionenkörper-Analogon}

Über Funktionenkörpern $k(C)$ (wobei $C$ eine Kurve über einem endlichen Körper
$\mathbb{F}_q$ ist) ist die $abc$-Vermutung ein Theorem, das aus dem Satz von
Mason--Stothers folgt (Mason 1983~\cite{Mason1984}):

\begin{theorem}[Satz von Mason--Stothers]
  Sei $k$ ein Körper der Charakteristik $0$. Wenn $a,b,c\in k[t]$ teilerfremde Polynome
  mit $a+b=c$ und $\max(\deg a,\deg b,\deg c)>0$ sind, dann gilt
  \[
    \max(\deg a,\deg b,\deg c) \;\leq\; \deg\operatorname{rad}(abc) - 1.
  \]
\end{theorem}

Dieser Satz ist deutlich stärker als die $abc$-Vermutung über Zahlkörpern (der Exponent
ist genau $1$, nicht $1+\varepsilon$) und wird durch ein einfaches Wronskian-Argument
bewiesen. Die Analogie legt nahe, dass ein arithmetisches Analogon des Wronskians der
Schlüssel zum Zahlkörper-Fall sein könnte.

% ================================================================
\section{Die $abc$-Vermutung und das Langlands-Programm}
% ================================================================

Die $abc$-Vermutung kann in der Sprache der Höhen auf der projektiven Geraden
$\mathbb{P}^1$ minus drei Punkte formuliert werden, was auf den Kontext von
\emph{Hyperbolischem Mordell} oder \emph{Shafarevich--Faltings} verweist.
Verbindungen umfassen:

\begin{itemize}[nosep]
  \item \textbf{Vojtas Vermutung:} Eine Verallgemeinerung von $abc$ auf rationale Punkte
    höherdimensionaler Varietäten; impliziert Faltings' Satz und mehr.
  \item \textbf{$S$-Einheitsgleichung:} Über Zahlkörpern ist die Menge
    $\{a+b=c : a,b,c\in\mathcal{O}_S^*\}$ endlich (Thue--Mahler--$S$-Einheitssatz);
    $abc$ liefert quantitative Schranken.
  \item \textbf{Weilsche Höhenmaschine:} Die Höhe $h(a,b,c)$ ist eine Weil-Höhe auf
    $\mathbb{P}^1$; Vojtas Vermutung in Dimension $1$ entspricht genau $abc$.
\end{itemize}

% ================================================================
\section{Aktueller Beweisstand}
% ================================================================

\begin{itemize}[nosep]
  \item \textbf{Bewiesen:} Schwache $abc$ (Exponent $2$, via Baker), Mason--Stothers
    (Funktionenkörper), Szpiro für CM-Kurven (Goldfeld--Szpiro), Fermat $n\geq 3$
    (Wiles, unbedingt).
  \item \textbf{Bedingt auf $abc$:} FLT für $n\geq 6$ (viel einfacher), Wieferich-Primzahlen
    (endlich viele), Halls Vermutung, verallgemeinerter Catalan.
  \item \textbf{Mochizukis IUT (2012):} Behaupteter Beweis; in PRIMS 2021 veröffentlicht;
    Konsens nicht erreicht; Einwand von Scholze--Stix ungeklärt.
  \item \textbf{Offen:} Die $abc$-Vermutung gilt als nicht bewiesen.
\end{itemize}

% ================================================================
\section{Schlussfolgerung}
% ================================================================

Die $abc$-Vermutung nimmt eine singuläre Position in der modernen Zahlentheorie ein:
eine einzige Ungleichung mit enormer Reichweite, die additive und multiplikative Strukturen
der ganzen Zahlen auf eine Weise verbindet, die sowohl elementar zu formulieren als auch
außerordentlich tief ist. Ihr Beweis würde weite Teile der diophantischen Geometrie neu
ordnen. Ob Mochizukis IUT letztendlich einen breiten Konsens erreicht oder ein neuer Ansatz
entsteht, bleibt abzuwarten.

\begin{thebibliography}{99}

\bibitem{Oesterle1988}
J.~Oesterlé,
\textit{Nouvelles approches du théorème de Fermat},
Séminaire Bourbaki, Exposé \textbf{694} (1988), 165--186.

\bibitem{Masser1985}
D.~W.~Masser,
\textit{Open problems},
in \textit{Proc.\ Sympos.\ Analytic Number Theory} (W.~W.~L.~Chen, Hrsg.),
Imperial College, London, 1985.

\bibitem{Granville1990}
A.~Granville,
\textit{The Beal conjecture and the $abc$ conjecture},
unveröffentlichte Notiz, 1990.
Siehe auch: A.~Granville und T.~J.~Tucker,
\textit{It's as easy as $abc$},
Notices AMS \textbf{49} (2002), 1224--1231.

\bibitem{Szpiro1981}
L.~Szpiro,
\textit{Séminaire sur les pinceaux de courbes de genre au moins deux},
Astérisque \textbf{86} (1981).

\bibitem{Silverman1994}
J.~H.~Silverman,
\textit{Advanced Topics in the Arithmetic of Elliptic Curves},
Springer, GTM \textbf{151}, 1994.

\bibitem{Nitaj1996}
A.~Nitaj,
\textit{La conjecture $abc$},
Enseign.\ Math.\ \textbf{42} (1996), 3--24.

\bibitem{Mason1984}
R.~C.~Mason,
\textit{Diophantine Equations over Function Fields},
LMS Lecture Note Series \textbf{96}, Cambridge University Press, 1984.

\bibitem{Mochizuki2012a}
S.~Mochizuki,
\textit{Inter-Universal Teichmüller Theory I: Construction of Hodge Theaters},
Preprint (2012), RIMS Kyoto.

\bibitem{Mochizuki2012b}
S.~Mochizuki,
\textit{Inter-Universal Teichmüller Theory II: Hodge-Arakelov-Theoretic Evaluation},
Preprint (2012), RIMS Kyoto.

\bibitem{Mochizuki2012c}
S.~Mochizuki,
\textit{Inter-Universal Teichmüller Theory III: Canonical Splittings of the Log-Theta-Lattice},
Preprint (2012), RIMS Kyoto.

\bibitem{Mochizuki2012d}
S.~Mochizuki,
\textit{Inter-Universal Teichmüller Theory IV: Log-Volume Computations and Set-Theoretic Foundations},
Preprint (2012), RIMS Kyoto.

\bibitem{ScholzeStix2018}
P.~Scholze und J.~Stix,
\textit{Why abc is still a conjecture},
Bericht (2018), verfügbar unter \url{https://www.math.uni-bonn.de/people/scholze/WhyABCisStillaConjecture.pdf}.

\bibitem{Mochizuki2018rebuttal}
S.~Mochizuki,
\textit{Comments on the manuscript by Scholze--Stix},
Bericht (2018), verfügbar auf der RIMS-Kyoto-Website.

\bibitem{Tan2022}
F.~Tan,
\textit{A remark on Mochizuki's Corollary 3.12},
Preprint, 2022.

\bibitem{Vojta1987}
P.~Vojta,
\textit{Diophantine Approximations and Value Distribution Theory},
Lecture Notes in Mathematics \textbf{1239}, Springer, 1987.

\end{thebibliography}

\end{document}
